% SEGMENT OLP-0676-S01
% Part: proof-theory
% Chapter: normalization
% Section: reductions

\documentclass[../../../include/open-logic-section]{subfiles}

\begin{document}

\olfileid{pt}{nor}{red}

\olsection{குறைப்பு உருமாற்றங்கள்}

\Log{N1i} இல் உள்ள !!a{proof} இன் வெட்டு
ஒன்றின் நீளம்~$1$ என்றால், அது ஒரே !!{formula}
நிகழ்வைக் கொண்ட தொடர். அந்த !!{formula} ஓர்
!!a{elimination} அனுமானத்தின் முதன்முன்கோள்.
முதலில் பரிசீலிக்கும் நேர்வில் அது ஓர்
!!a{introduction} அனுமானத்தின் முடிவும் ஆகும்;
பொய்யிலிருந்து வருவிக்கும் நேர்வை இறுதியில் பார்ப்போம்.
இயல்பாக்க நிறுவலில் இத்தகைய !!{introduction}/
!!{elimination} இணையில் முடியும் துணை-!!{proof}
ஐ, அவ்விரு அனுமானங்களும் நீக்கப்பட்ட புதிய
!!{proof} ஆல் மாற்றிக் குறைக்கிறோம். இந்த மாற்றத்தால்
புதிய வெட்டுகள் உருவாகலாம். ஆனால் அவற்றின் நிலை,
குறைக்கப்படும் வெட்டின் நிலையைவிடத் தாழ்ந்தது என்பதை
கீழே உறுதிசெய்வோம். எனவே உரிய மிகுநிலை வெட்டைத்
தேர்ந்தெடுத்தால், புதிய !!{proof} இன் வெட்டு நீளம்
குறையும் (நீளம்~$1$ கொண்ட மிகுநிலை வெட்டு
நீக்கப்படுவதால்); அல்லது, அந்த வெட்டே ஒரே மிகுநிலை
வெட்டாக இருந்தால், !!{proof} இன் வெட்டு நிலை குறையும்.

% SEGMENT OLP-0676-S02
எடுத்துக்காட்டாக, தொடர்புடைய அனுமானங்கள்
\Intro{\lif}, \Elim{\lif} ஆகவும்,
$!A \ident !B \lif !C$ ஆகவும், துணை-!!{proof}~$\delta_1$
பின்வருமாறு முடிவதாகவும் கொள்வோம்:
\[
\AxiomC{$\Discharge{!B}{x}$}
\RightLabel{$\delta_2$}
\DeduceC{$!C$}
\DischargeRule{\Intro{\lif}}{x}
\UnaryInfC{$!B \lif !C$}
\AxiomC{}
\RightLabel{$\delta_3$}
\DeduceC{$!B$}
\RightLabel{\Elim{\lif}}
\BinaryInfC{$!C$}
\DisplayProof
\]
இந்த வெட்டை அகற்ற, $\delta$ இல் உள்ள
துணைநிறுவல்~$\delta_1$ ஐப் பின்வருவதாக மாற்றுகிறோம்:
\[
\AxiomC{}
\RightLabel{$\delta_3$}
\DeduceC{$!B$}
\RightLabel{$\Subst{\delta_2}{\delta_3}{!B^x}$}
\DeduceC{$!C$}
\DisplayProof
\]
$\Subst{\delta_2}{\delta_3}{!B^x}$ என்பதன் பொருள்:
$\delta_2$ இல் $x$ என்று குறிக்கப்பட்ட ஒவ்வொரு
$!B$ எடுகோளுக்குப் பதிலாகவும் முழு !!{proof}~$\delta_3$
ஐ வைத்துப் பெறும் !!{proof}. இந்தப் புதிய !!{proof} இல்
$x$ குறியுடைய திறந்த $!B$ எடுகோள்கள் எஞ்சாது.
ஒரே குறியால் இரண்டு வெவ்வேறு !!{formula}s எடுகோள்களாக
வராது என்று வழக்கம்போலக் கொள்வதால், வேறு எந்த
எடுகோளுக்கும் $x$ குறி இருக்காது. ஆகவே புதிய !!{proof}
இன் திறந்த எடுகோள்கள்~$\delta_1$ இன்வையே.
$\delta_2$ இல் எடுகோள்களை விடுவிக்கும் ஒவ்வோர் அனுமானமும்
$\Subst{\delta_2}{\delta_3}{!B^x}$ இலும் அதே
எடுகோள்களை விடுவிக்கிறது; $\delta_3$ இல் எடுகோள்களை
விடுவிக்கும் ஒவ்வோர் அனுமானமும்
$\Subst{\delta_2}{\delta_3}{!B^x}$ இல் தோன்றும்
ஒன்று அல்லது பல நகல்களில் தொடர்புடைய எடுகோள்களை
விடுவிக்கிறது. ஒட்டப்பட்ட நகல்களில் தனிமாறிகள்
மோதுமானால், அவற்றுக்கு வேறுபட்ட புதுப் பெயர்களை
ஒதுக்குகிறோம்.

% SEGMENT OLP-0676-S03
$\delta_2$ அல்லது $\delta_3$ இன் ஒரு நகலுக்குள்
முழுவதுமாகவோ, அல்லது $\delta$ இல்~$\delta_1$
க்கு முற்றிலும் வெளியிலோ இருக்கும் ஒவ்வொரு வெட்டும்
தனித்தனியாக மாறாது: அதே !!{formula}s மற்றும்
அனுமானங்களைக் கொண்டிருப்பதால் $\delta$ இன்
ஒத்த வெட்டுக்குரிய அதே நிலையும் நீளமும் கொண்டது.
ஆனால் ஒட்டப்படும் நிறுவலின் நகல்களின் எண்ணிக்கை மாறலாம்;
அதனால் மொத்த வெட்டு நீளம் குறைந்தது என்று இந்தக்
குறிப்பிலிருந்து மட்டும் பின்தொடராது. நீளம்~$1$
கொண்ட ஒரு மிகுநிலை வெட்டை நீக்கியுள்ளோம். $\delta$
இன் வெட்டு நீளம்~$1$ ஆகவும் இதுவே ஒரே மிகுநிலை
வெட்டாகவும் இருந்தால் வெட்டு நிலை குறையும்; மற்ற
நேர்வுகளில் நகலாக்கத்தின் விளைவையும் கணக்கில் கொள்ள
வேண்டும்.

% SEGMENT OLP-0676-S04
ஆயினும் இரு நிகழ்வுகளைக் கவனிக்க வேண்டும்:
\begin{enumerate}
  \item $!B^x$ என்ற எல்லா எடுகோள்களுக்கும்
  $\delta_3$ ஐப் பதிலிடும்போது, $\delta_3$ இல் உள்ள
  வெட்டுகளும் நகலாகின்றன. அவற்றில் மிகுநிலை
  வெட்டுகள் இருந்தால், புதிய !!{proof} இன் வெட்டு
  நீளம் பழைய !!{proof} இன் வெட்டு நீளத்தைவிட
  அதிகமாகலாம்.
  இதைத் தவிர்க்க, \emph{மிக வலப்புறமான} மிகுநிலை
  வெட்டுகளுக்கு மட்டுமே குறைப்பு உருமாற்றத்தைப்
  பயன்படுத்துகிறோம். அதாவது குறைக்கப்படும் வெட்டுக்கு
  வலப்புறமாக $\delta$ இன் எந்தக் கிளையிலும் மற்றொரு
  மிகுநிலை வெட்டு இருக்கக்கூடாது. $\delta_3$ இல்
  மிகுநிலை வெட்டு இருந்தால் அது இங்கு குறைக்கப்படும்
  வெட்டின் வலப்புறம் இருக்கும்; எனவே தேர்ந்தெடுக்கப்பட்ட
  வெட்டு மிக வலப்புறமானதாக இருக்காது. இவ்வாறு
  தேர்ந்தெடுப்பதால் !!{proof} இன் வெட்டு நீளத்தைக்
  குறைக்கவோ, வெட்டு நிலையையே குறைக்கவோ முடிகிறது.
  \item $\delta_2$ இல் $!B^x$ எடுகோள் ஓர்
  !!a{elimination} அனுமானத்தின் முதன்முன்கோளாக
  இருந்து, $\delta_3$ இன் கடைசி அனுமானம்
  \Elim{\lor}, \Elim{\lexists} அல்லது ஓர்
  !!a{introduction} அனுமானமாக இருந்தால்,
  அந்த இடத்தில்~$\delta_3$ ஐ $\delta_2$ உடன்
  ஒட்டுவதால் புதிய வெட்டு தோன்றும். $\delta_2$
  இல் வெறும் எடுகோளாக இருந்தது இப்போது
  நீளம்~$1$ கொண்ட வெட்டாகலாம் (ஓர்
  !!a{introduction} அனுமானத்தின் முடிவும்
  !!a{elimination} அனுமானத்தின் முதன்முன்கோளும்);
  அல்லது வெட்டுத் தொடரின் கடைசி !!{formula}
  நிகழ்வாகலாம் (ஓர் !!a{elimination} அனுமானத்தின்
  முதன்முன்கோளும் \Elim{\lor} அல்லது
  \Elim{\lexists} இன் முடிவும்). $!B^x$ என்பது
  \Elim{\lor} அல்லது \Elim{\lexists} இன்
  சிறுமுன்கோளாக இருந்தால், அது ஏற்கெனவே ஒரு
  தொடரின் முதல் !!{formula} நிகழ்வாகவும்,
  ஒருவேளை வெட்டுத் தொடரின் முதல் நிகழ்வாகவும்
  இருந்தது. புதிய !!{proof} இல் அந்தத் தொடர்
  நீளமாகலாம். ஆனால் புதிதாக உண்டான வெட்டின்
  அல்லது நீண்ட தொடரின் !!{formula}~$!B$ ஆகும்;
  மேலும் $\depth{!B} < \depth{!B \lif !C}$.
  எனவே வெட்டுகள் சேர்ந்தாலும் அவற்றின் நீளம்
  அதிகரித்தாலும், புதிதாக வந்தவை எதுவும்
  மிகுநிலை வெட்டுகள் அல்ல. ஆகவே வெட்டு நிலை
  குறையும்; அதே நிலையில் வெட்டுகள் எஞ்சினால்,
  வெட்டு நீளம் குறையும்.
\end{enumerate}

% SEGMENT OLP-0676-S05
நீளம்~$1$ கொண்ட வெட்டைக் குறைக்கும் செயல்
\emph{குறைப்பு உருமாற்றம்} (அல்லது
\emph{சுற்றுவழி உருமாற்றம்}) எனப்படுகிறது.
சில விதி இணைகளுக்கான வடிவங்கள்
\olref{tab:reduction-conversions} இல் தரப்பட்டுள்ளன.
(மீதியுள்ளவைப் பயிற்சிகளாக விடப்பட்டுள்ளன.)

\begin{table}
\begin{multline*}
\bottomAlignProof
\AxiomC{$\Discharge{!B}{x}$}
\RightLabel{$\delta_2$}
\shortDeduce
\DeduceC{$!C$}
\DischargeRule{\Intro{\lif}}{x}
\UnaryInfC{$!B \lif !C$}
\AxiomC{}
\RightLabel{$\delta_3$}
\shortDeduce
\DeduceC{$!B$}
\RightLabel{\Elim{\lif}}
\BinaryInfC{$!C$}
\DisplayProof
\quad \longrightarrow\quad
\shoveright{\bottomAlignProof
\AxiomC{}
\RightLabel{$\delta_3$}
\shortDeduce
\DeduceC{$!B$}
\RightLabel{$\Subst{\delta_2}{\delta_3}{!B^x}$}
\shortDeduce
\DeduceC{$!C$}
\DisplayProof}
\\[4ex]
\shoveleft{\bottomAlignProof
\AxiomC{}
\RightLabel{$\delta_2$}
\shortDeduce
\DeduceC{$!B$}
\AxiomC{}
\RightLabel{$\delta_3$}
\shortDeduce
\DeduceC{$!C$}
\RightLabel{\Intro{\land}}
\BinaryInfC{$!B \land !C$}
\RightLabel{\Elim{\land}[1]}
\UnaryInfC{$!B$}
\DisplayProof
\quad \longrightarrow\quad
\bottomAlignProof
\AxiomC{}
\RightLabel{$\delta_2$}
\shortDeduce
\DeduceC{$!B$}
\DisplayProof}\qquad
\shoveright{\bottomAlignProof
\AxiomC{}
\RightLabel{$\delta_2$}
\shortDeduce
\DeduceC{$!B$}
\AxiomC{}
\RightLabel{$\delta_3$}
\shortDeduce
\DeduceC{$!C$}
\RightLabel{\Intro{\land}}
\BinaryInfC{$!B \land !C$}
\RightLabel{\Elim{\land}[2]}
\UnaryInfC{$!C$}
\DisplayProof
\quad \longrightarrow\quad
\bottomAlignProof
\AxiomC{}
\RightLabel{$\delta_3$}
\shortDeduce
\DeduceC{$!C$}
\DisplayProof}
\\[4ex]
\bottomAlignProof
\AxiomC{}
\RightLabel{$\delta_2$}
\shortDeduce
\DeduceC{$!B$}
\RightLabel{\Intro{\lor}[1]}
\UnaryInfC{$!B \lor !C$}
\AxiomC{$\Discharge{!B}{x}$}
\RightLabel{$\delta_3$}
\shortDeduce
\DeduceC{$!D$}
\AxiomC{$\Discharge{!C}{y}$}
\RightLabel{$\delta_4$}
\shortDeduce
\DeduceC{$!D$}
\DischargeRule{\Elim{\lor}}{x\,y}
\TrinaryInfC{$!D$}
\DisplayProof
\quad \longrightarrow\quad
\shoveright{\bottomAlignProof
\AxiomC{}
\RightLabel{$\delta_2$}
\shortDeduce
\DeduceC{$!B$}
\RightLabel{$\Subst{\delta_3}{\delta_2}{!B^x}$}
\shortDeduce
\DeduceC{$!D$}
\DisplayProof}
\\[2ex]
\shoveleft{\bottomAlignProof
\AxiomC{}
\RightLabel{$\delta_2$}
\shortDeduce
\DeduceC{$!B(c)$}
\RightLabel{\Intro{\lforall}}
\UnaryInfC{$\lforall[x][!B(x)]$}
\RightLabel{\Elim{\lforall}}
\UnaryInfC{$!B(t)$}
\DisplayProof
\quad \longrightarrow\quad
\bottomAlignProof
\AxiomC{}
\RightLabel{$\Subst{\delta_2}{t}{c}$}
\shortDeduce
\DeduceC{$!B(t)$}
\DisplayProof}
\\[4ex]
\shoveright{\bottomAlignProof
\AxiomC{}
\RightLabel{$\delta_2$}
\shortDeduce
\DeduceC{$!B(t)$}
\RightLabel{\Intro{\lexists}}
\UnaryInfC{$\lexists[x][!B(x)]$}
\AxiomC{$\Discharge{!B(c)}{x}$}
\RightLabel{$\delta_3$}
\shortDeduce
\DeduceC{$!D$}
\DischargeRule{\Elim{\lexists}}{x}
\BinaryInfC{$!D$}
\DisplayProof
\quad \longrightarrow\quad
\bottomAlignProof
\AxiomC{}
\RightLabel{$\delta_2$}
\shortDeduce
\DeduceC{$!B(t)$}
\RightLabel{$\Subst{\Subst{\delta_3}{t}{c}}{\delta_2}{!B(t)^x}$}
\shortDeduce
\DeduceC{$!D$}
\DisplayProof}
\end{multline*}
\caption{\Log{N1i} இன் சில குறைப்பு உருமாற்றங்கள்.}
\ollabel{tab:reduction-conversions}
\end{table}

% SEGMENT OLP-0676-S06
\begin{prob}
\Intro{\lnot} ஐத் தொடர்ந்து \Elim{\lnot}
வரும்போது பொருந்தும் குறைப்பு உருமாற்றத்தைக்
கொடுக்கவும்.
\end{prob}

இன்னும் ஒரு நேர்வைப் பரிசீலிக்க வேண்டும்.
வெட்டின் வரையறை, நீளம்~$1$ ஆகவும் $!A$
என்பது \FalseInt{} இன் முடிவாகவும் ஓர்
!!a{elimination} அனுமானத்தின் முதன்முன்கோளாகவும்
இருக்கும் நேர்வையும் உள்ளடக்குகிறது. $!A$
ஓர் !!a{introduction} விதியின் முடிவாக இருக்கும்
நேர்வைப் போலவே இதையும் கையாளலாம். எடுத்துக்காட்டாக,
பின்வருவதை மாற்றுக:
\[
\bottomAlignProof
\AxiomC{}
\RightLabel{$\delta_2$}
\DeduceC{$\lfalse$}
\UnaryInfC{$!B \lif !C$}
\AxiomC{}
\RightLabel{$\delta_3$}
\DeduceC{$!B$}
\RightLabel{\Elim{\lif}}
\BinaryInfC{$!C$}
\DisplayProof
\quad\text{இதற்கு பதில்}\quad
\bottomAlignProof
\AxiomC{}
\RightLabel{$\delta_2$}
\DeduceC{$\lfalse$}
\UnaryInfC{$!C$}
\DisplayProof
\]

% SEGMENT OLP-0676-S07
ஆனால் $!A \ident (!B \lor !C)$ அல்லது
$!A \ident \lexists[x][!B(x)]$ என்ற
நேர்வுகளில் சிறிது கவனம் தேவை. எடுத்துக்காட்டாக,
பின்வருவதை மாற்றினால்
\[
\bottomAlignProof
\AxiomC{}
\RightLabel{$\delta_2$}
\DeduceC{$\lfalse$}
\RightLabel{\FalseInt}
\UnaryInfC{$!B \lor !C$}
\AxiomC{$\Discharge{!B}{x}$}
\RightLabel{$\delta_3$}
\DeduceC{$!D$}
\AxiomC{$\Discharge{!C}{y}$}
\RightLabel{$\delta_4$}
\DeduceC{$!D$}
\DischargeRule{\Elim{\lor}}{x\,y}
\TrinaryInfC{$!D$}
\DisplayProof
\quad\text{இதற்கு பதில்}\quad
\bottomAlignProof
\AxiomC{}
\RightLabel{$\delta_2$}
\DeduceC{$\lfalse$}
\RightLabel{\FalseInt}
\UnaryInfC{$!D$}
\DisplayProof
\]
வெட்டு !!{formula}~$!D$ கொண்ட புதிய வெட்டு
உண்டாகலாம்; $\depth{!D} < \depth{!B \lor !C}$!{}
என்பதற்கு உத்தரவாதம் இல்லை. ஆகவே இங்கேயும்
\Intro{\lor}/\Elim{\lor} குறைப்பு நேர்வைப் போலச்
செயல்பட்டு, பின்வருவதாக மாற்ற வேண்டும்:
\[
\bottomAlignProof
\AxiomC{}
\RightLabel{$\delta_2$}
\DeduceC{$\lfalse$}
\RightLabel{\FalseInt}
\UnaryInfC{$!B$}
\RightLabel{$\delta_3$}
\DeduceC{$!D$}
\DisplayProof
\quad\text{அல்லது}\quad
\bottomAlignProof
\AxiomC{}
\RightLabel{$\delta_2$}
\DeduceC{$\lfalse$}
\RightLabel{\FalseInt}
\UnaryInfC{$!C$}
\RightLabel{$\delta_4$}
\DeduceC{$!D$}
\DisplayProof
\]

\end{document}
